% Bagian: computability
% Bab: recursive-functions
% Subbagian: trees

\documentclass[../../../include/open-logic-section]{subfiles}

\begin{document}

\olfileid[id]{cmp}{rec}{tre}
\olsection{Pohon}

Kadang-kadang berguna untuk merepresentasikan pohon sebagai bilangan asli,
seperti halnya kita dapat merepresentasikan barisan dengan bilangan serta
merepresentasikan sifat dan operasi pada barisan melalui relasi dan fungsi
rekursif primitif pada kode-kodenya. Kita akan menggunakan barisan dan
kode-kodenya untuk melakukan hal ini. Suatu pohon dapat berupa satu simpul
(mungkin dengan label), atau berupa suatu simpul (mungkin dengan label) yang
terhubung dengan sejumlah subpohon. Simpul tersebut disebut \emph{akar} pohon,
dan subpohon yang terhubung dengannya disebut \emph{subpohon langsung}nya.

Kita mengodekan pohon secara rekursif sebagai barisan
$\tuple{k,d_1,\dots,d_k}$, dengan $k$ banyaknya subpohon langsung dan
$d_1$, \dots,~$d_k$ kode-kode subpohon langsung tersebut. Jika simpul-simpul
mempunyai label, label dapat dicantumkan setelah subpohon langsung. Jadi,
pohon yang hanya terdiri atas satu simpul dengan label~$l$ akan dikodekan oleh
$\tuple{0,l}$, dan pohon yang terdiri atas suatu akar (berlabel~$l_1$) yang
terhubung dengan dua simpul tunggal (berlabel $l_2$, $l_3$) akan dikodekan
oleh $\tuple{2,\tuple{0,l_2},\tuple{0,l_3},l_1}$.

\begin{prop}
  \ollabel{prop:subtreeseq}
  Fungsi $\fn{SubtreeSeq}(t)$, yang mengembalikan kode suatu barisan yang
  elemen-elemennya merupakan kode semua subpohon dari pohon berkode~$t$,
  bersifat rekursif primitif.
\end{prop}

\begin{proof}
  Pertama, perhatikan bahwa
  $\fn{ISubtrees}(t)=\fn{subseq}(t,1,(t)_0)$ bersifat rekursif primitif dan
  mengembalikan kode-kode subpohon langsung dari suatu pohon~$t$. Sekarang
  kita dapat mendefinisikan fungsi pembantu $\fn{hSubtreeSeq}(t,n)$ yang
  menghitung barisan semua subpohon yang berjarak paling jauh $n$ simpul dari
  akar. Barisan subpohon dari~$t$ yang berjarak paling jauh $0$ simpul dari
  akar---dengan kata lain, yang dimulai pada akar~$t$---merupakan barisan yang
  hanya terdiri atas~$t$. Untuk memperoleh barisan semua subpohon sampai
  tingkat~$n+1$ dari~$t$, kita menggabungkan subpohon sampai tingkat~$n$
  dengan barisan yang terdiri atas semua subpohon langsung dari subpohon
  sampai tingkat~$n$. Untuk memperoleh daftar semua subpohon langsung ini,
  perhatikan bahwa jika $f(x)$ merupakan fungsi rekursif primitif yang
  mengembalikan kode barisan, maka
  $g_f(s,k)=f((s)_0)\concat\dots\concat f((s)_{k-1})$ untuk $k>0$ juga
    bersifat rekursif primitif:
    \begin{align*}
      g_f(s,0) & = 0\\
      g_f(s,k+1) & = g_f(s,k) \concat f((s)_k).
    \end{align*}
    Sebagai contoh, jika $s$ merupakan barisan pohon, maka
    $h(s)=g_{\fn{ISubtrees}}(s,\len{s})$ menghasilkan barisan semua subpohon
    langsung dari elemen-elemen~$s$. Kita dapat menggunakannya untuk
    mendefinisikan $\fn{hSubtreeSeq}$ dengan
    \begin{align*}
      \fn{hSubtreeSeq}(t, 0) & = \tuple{t} \\
      \fn{hSubtreeSeq}(t, n+1) & = \fn{hSubtreeSeq}(t, n) \concat
      h(\fn{hSubtreeSeq}(t, n)).
    \end{align*}
    Tingkat maksimum subpohon dalam pohon yang dikodekan oleh~$t$, yakni jarak
    maksimum antara akar dan suatu simpul daun, dibatasi oleh kode~$t$. Jadi,
    barisan kode semua subpohon dari pohon berkode~$t$ diberikan oleh
    $\fn{hSubtreeSeq}(t,t)$.
\end{proof}

\begin{prob}
  Definisi $\fn{hSubtreeSeq}$ dalam bukti
  \olref[cmp][rec][tre]{prop:subtreeseq} pada umumnya memuat pengulangan.
  Berikan definisi lain yang menjamin bahwa kode suatu subpohon hanya muncul
  satu kali dalam daftar yang dihasilkan.
\end{prob}

\end{document}
